\documentclass[12pt,a4paper]{article}

% ==================== PACKAGES ====================
\usepackage[utf8]{inputenc}
\usepackage[vietnamese]{babel}
\usepackage{amsmath,amssymb,amsthm}
\usepackage{geometry}
\usepackage{enumitem}
\usepackage[most]{tcolorbox}
\usepackage{hyperref}
\usepackage{fancyhdr}
\usepackage{graphicx}
\usepackage{tikz}
\usepackage{eso-pic}
\usepackage{xcolor}

\geometry{a4paper, margin=2cm, bottom=2.5cm}
\hypersetup{colorlinks=true, linkcolor=blue!70!black, urlcolor=blue}
\setlist[itemize]{topsep=4pt, itemsep=2pt}
\setlist[enumerate]{topsep=4pt, itemsep=2pt}

% ==================== WATERMARK ====================
\AddToShipoutPictureFG{%
  \AtPageCenter{%
    \begin{tikzpicture}[remember picture, overlay]
      \node[rotate=45, scale=1, opacity=0.06]{%
        \fontsize{60}{72}\selectfont\bfseries\sffamily Đặng Tấn Quí%
      };
    \end{tikzpicture}%
  }%
}

% ==================== HEADER / FOOTER ====================
\pagestyle{fancy}
\fancyhf{}
\fancyhead[L]{\small\textit{Toán cao cấp 1 -- Đại số tuyến tính}}
\fancyhead[R]{\small\textit{Đặng Tấn Quí}}
\fancyfoot[C]{\thepage}
\fancyfoot[L]{\footnotesize\textcopyright\ Đặng Tấn Quí}
\fancyfoot[R]{\footnotesize SĐT: 0377\,122\,210}
\renewcommand{\headrulewidth}{0.4pt}
\renewcommand{\footrulewidth}{0.4pt}

\fancypagestyle{plain}{%
  \fancyhf{}
  \fancyfoot[C]{\thepage}
  \fancyfoot[L]{\footnotesize\textcopyright\ Đặng Tấn Quí}
  \fancyfoot[R]{\footnotesize SĐT: 0377\,122\,210}
  \renewcommand{\headrulewidth}{0pt}
  \renewcommand{\footrulewidth}{0.4pt}
}

% ==================== CUSTOM ENVIRONMENTS ====================
\newtcolorbox{dinhnghia}[1][]{%
  enhanced, breakable,
  colback=blue!3!white, colframe=blue!60!black,
  fonttitle=\bfseries, title={Định nghĩa}, #1%
}

\newtcolorbox{dinhly}[1][]{%
  enhanced, breakable,
  colback=green!3!white, colframe=green!50!black,
  fonttitle=\bfseries, title={Định lý}, #1%
}

\newtcolorbox{tinhchat}[1][]{%
  enhanced, breakable,
  colback=yellow!5!white, colframe=orange!50!black,
  fonttitle=\bfseries, title={Tính chất}, #1%
}

\newtcolorbox{phuongphap}[1][]{%
  enhanced, breakable,
  colback=gray!5!white, colframe=gray!50!black,
  fonttitle=\bfseries, title={Phương pháp}, #1%
}

\newtcolorbox{luuy}[1][]{%
  enhanced, breakable,
  colback=red!3!white, colframe=red!50!black,
  fonttitle=\bfseries, title={Lưu ý}, #1%
}

\newtcolorbox{congthuc}[1][]{%
  enhanced, breakable,
  colback=cyan!3!white, colframe=cyan!50!black,
  fonttitle=\bfseries, title={Công thức}, #1%
}

\newtcolorbox{vidu}[1][]{%
  enhanced, breakable,
  colback=violet!3!white, colframe=violet!50!black,
  fonttitle=\bfseries, title={Ví dụ}, #1%
}

% ==================== DOCUMENT ====================
\begin{document}

% ==================== TRANG BÌA ====================
\thispagestyle{empty}
\begin{tikzpicture}[remember picture, overlay]
  \fill[blue!80!black] (current page.south west) rectangle (current page.north east);
  \fill[blue!60!cyan, opacity=0.8]
    (current page.south west) -- (current page.north west) --
    ([xshift=-3cm]current page.north east) -- (current page.south east) -- cycle;
  \foreach \r/\o in {6/0.05, 4.5/0.08, 3/0.12} {
    \fill[white, opacity=\o] ([xshift=2cm, yshift=-2cm]current page.north east) circle (\r cm);
  }
  \draw[white, line width=2pt] ([yshift=-5cm]current page.north west) ++(2cm,0) -- ++(17cm,0);
  \draw[white, line width=2pt] ([yshift=-16cm]current page.north west) ++(2cm,0) -- ++(17cm,0);
  \node[anchor=west, white, font=\fontsize{28}{34}\bfseries\selectfont]
    at ([xshift=3cm, yshift=-7.5cm]current page.north west)
    {TOÁN CAO CẤP 1};
  \node[anchor=west, white, font=\fontsize{20}{26}\selectfont]
    at ([xshift=3cm, yshift=-9cm]current page.north west)
    {Đại số tuyến tính --- Năm 1, HK1};
  \node[anchor=west, white, font=\fontsize{18}{24}\selectfont]
    at ([xshift=3cm, yshift=-10cm]current page.north west)
    {Tài liệu lý thuyết};
  \node[anchor=west, white, font=\Large\bfseries]
    at ([xshift=3cm, yshift=-13cm]current page.north west) {Biên soạn:};
  \node[anchor=west, yellow!80!white, font=\fontsize{24}{30}\bfseries\selectfont]
    at ([xshift=3cm, yshift=-14.5cm]current page.north west) {ĐẶNG TẤN QUÍ};
  \node[anchor=west, white!90, font=\large]
    at ([xshift=3cm, yshift=-18cm]current page.north west) {SĐT: 0377\,122\,210};
  \node[anchor=south, white!80, font=\small]
    at ([yshift=1.5cm]current page.south) {Tài liệu có bản quyền --- Cấm sao chép khi chưa được sự đồng ý của tác giả};
  \node[anchor=south east, white, font=\Large\bfseries]
    at ([xshift=-2cm, yshift=3cm]current page.south east) {\the\year};
\end{tikzpicture}

\newpage
\tableofcontents
\newpage

%% ================================================================
%%  CHƯƠNG 1: MA TRẬN
%% ================================================================
\section{Ma trận}

\subsection{Định nghĩa ma trận}

\begin{dinhnghia}
  \textbf{Ma trận} cỡ $m \times n$ (hoặc cấp $m \times n$) là bảng số hình chữ nhật gồm $m$ dòng và $n$ cột:
  \[
    A = \begin{pmatrix}
      a_{11} & a_{12} & \cdots & a_{1n} \\
      a_{21} & a_{22} & \cdots & a_{2n} \\
      \vdots & \vdots & \ddots & \vdots \\
      a_{m1} & a_{m2} & \cdots & a_{mn}
    \end{pmatrix} = (a_{ij})_{m \times n}.
  \]
  Phần tử $a_{ij}$ nằm ở dòng $i$, cột $j$.
\end{dinhnghia}

\begin{dinhnghia}
  \textbf{Ma trận vuông} cấp $n$: $m = n$. Đường chéo chính: $a_{11}, a_{22}, \ldots, a_{nn}$.
\end{dinhnghia}

\begin{tinhchat}[title={Các loại ma trận đặc biệt}]
  \begin{itemize}
    \item \textbf{Ma trận không} $\mathbf{0}$: mọi phần tử bằng 0.
    \item \textbf{Ma trận đơn vị} $I_n$: ma trận vuông cấp $n$, đường chéo chính bằng 1, còn lại bằng 0.
    \item \textbf{Ma trận chuyển vị} $A^T$: $(A^T)_{ij} = a_{ji}$.
    \item \textbf{Ma trận đối xứng}: $A^T = A$.
  \end{itemize}
\end{tinhchat}

\subsection{Các phép toán trên ma trận}

\begin{dinhnghia}
  \textbf{Cộng ma trận} (cùng cỡ): $(A + B)_{ij} = a_{ij} + b_{ij}$.
\end{dinhnghia}

\begin{dinhnghia}
  \textbf{Nhân ma trận với số}: $(\lambda A)_{ij} = \lambda a_{ij}$.
\end{dinhnghia}

\begin{dinhnghia}
  \textbf{Nhân ma trận} $A_{m \times p}$ với $B_{p \times n}$:
  \[
    (AB)_{ij} = \sum_{k=1}^{p} a_{ik} b_{kj}.
  \]
  Điều kiện: số cột của $A$ = số dòng của $B$.
\end{dinhnghia}

\begin{tinhchat}
  \begin{itemize}
    \item $A(BC) = (AB)C$.
    \item $A(B + C) = AB + AC$.
    \item $(AB)^T = B^T A^T$.
    \item Nói chung $AB \ne BA$.
  \end{itemize}
\end{tinhchat}

%% ================================================================
%%  CHƯƠNG 2: ĐỊNH THỨC
%% ================================================================
\section{Định thức}

\subsection{Định nghĩa định thức}

\begin{dinhnghia}
  \textbf{Định thức} của ma trận vuông $A$ cấp $n$, ký hiệu $\det(A)$ hoặc $|A|$:
  \begin{itemize}
    \item Cấp 1: $\det(a) = a$.
    \item Cấp 2: $\det\begin{pmatrix} a & b \\ c & d \end{pmatrix} = ad - bc$.
    \item Cấp $n$: khai triển theo dòng $i$: $\det(A) = \sum_{j=1}^{n} (-1)^{i+j} a_{ij} M_{ij}$, với $M_{ij}$ là định thức con bù của $a_{ij}$.
  \end{itemize}
\end{dinhnghia}

\begin{dinhly}[title={Tính chất định thức}]
  \begin{itemize}
    \item $\det(A^T) = \det(A)$.
    \item Đổi chỗ hai dòng (hoặc hai cột) $\Rightarrow$ định thức đổi dấu.
    \item Nhân một dòng (cột) với $\lambda$ $\Rightarrow$ định thức nhân với $\lambda$.
    \item $\det(AB) = \det(A) \cdot \det(B)$.
    \item Ma trận có một dòng (cột) bằng 0 $\Rightarrow$ $\det(A) = 0$.
  \end{itemize}
\end{dinhly}

\begin{phuongphap}[title={Tính định thức}]
  \begin{enumerate}
    \item Khai triển theo dòng hoặc cột có nhiều số 0.
    \item Biến đổi sơ cấp đưa về ma trận tam giác (định thức = tích phần tử đường chéo).
  \end{enumerate}
\end{phuongphap}

%% ================================================================
%%  CHƯƠNG 3: HỆ PHƯƠNG TRÌNH TUYẾN TÍNH
%% ================================================================
\section{Hệ phương trình tuyến tính}

\subsection{Dạng ma trận}

\begin{dinhnghia}
  Hệ $m$ phương trình $n$ ẩn:
  \[
    \begin{cases}
      a_{11}x_1 + a_{12}x_2 + \cdots + a_{1n}x_n = b_1 \\
      a_{21}x_1 + a_{22}x_2 + \cdots + a_{2n}x_n = b_2 \\
      \vdots \\
      a_{m1}x_1 + a_{m2}x_2 + \cdots + a_{mn}x_n = b_m
    \end{cases}
    \quad \Leftrightarrow \quad Ax = b.
  \]
  Ma trận hệ số $A$, ma trận mở rộng $\overline{A} = [A \mid b]$.
\end{dinhnghia}

\begin{dinhly}[title={Định lý Kronecker--Capelli}]
  Hệ $Ax = b$ có nghiệm khi và chỉ khi $\mathrm{rank}(A) = \mathrm{rank}(\overline{A})$.
  \begin{itemize}
    \item Nếu $\mathrm{rank}(A) = \mathrm{rank}(\overline{A}) = n$: hệ có nghiệm duy nhất.
    \item Nếu $\mathrm{rank}(A) = \mathrm{rank}(\overline{A}) < n$: hệ có vô số nghiệm (số ẩn tự do = $n - \mathrm{rank}(A)$).
  \end{itemize}
\end{dinhly}

\begin{phuongphap}[title={Phương pháp Gauss}]
  Biến đổi sơ cấp trên dòng đưa ma trận mở rộng về dạng bậc thang, từ đó suy ra nghiệm.
\end{phuongphap}

\begin{dinhly}[title={Hệ Cramer}]
  Hệ $n$ phương trình $n$ ẩn $Ax = b$. Nếu $\det(A) \ne 0$ thì hệ có nghiệm duy nhất:
  \[
    x_j = \frac{\det(A_j)}{\det(A)},
  \]
  với $A_j$ là ma trận $A$ thay cột $j$ bởi cột $b$.
\end{dinhly}

%% ================================================================
%%  CHƯƠNG 4: KHÔNG GIAN VECTƠ VÀ ÁNH XẠ TUYẾN TÍNH
%% ================================================================
\section{Không gian vectơ}

\subsection{Định nghĩa và ví dụ}

\begin{dinhnghia}
  Tập $V$ với hai phép toán: cộng $(+)$ và nhân với vô hướng $(\cdot)$, gọi là \textbf{không gian vectơ} trên $\mathbb{R}$ nếu thỏa các tiên đề (tính đóng, giao hoán, kết hợp, phần tử không, phần tử đối, phân phối, đơn vị).
\end{dinhnghia}

\begin{vidu}
  $\mathbb{R}^n$, tập các ma trận $m \times n$, tập các đa thức bậc $\le n$.
\end{vidu}

\subsection{Độc lập tuyến tính, phụ thuộc tuyến tính}

\begin{dinhnghia}
  Hệ vectơ $\{v_1, \ldots, v_k\}$ \textbf{phụ thuộc tuyến tính} nếu tồn tại $\lambda_1, \ldots, \lambda_k$ không đồng thời bằng 0 sao cho $\lambda_1 v_1 + \cdots + \lambda_k v_k = 0$. Ngược lại là \textbf{độc lập tuyến tính}.
\end{dinhnghia}

\begin{dinhnghia}
  \textbf{Hạng} của hệ vectơ: số vectơ tối đa độc lập tuyến tính trong hệ. Ký hiệu $\mathrm{rank}\{v_1, \ldots, v_k\}$.
\end{dinhnghia}

\subsection{Cơ sở và số chiều}

\begin{dinhnghia}
  Hệ vectơ $\mathcal{B}$ là \textbf{cơ sở} của $V$ nếu $\mathcal{B}$ độc lập tuyến tính và mọi vectơ của $V$ biểu diễn được qua $\mathcal{B}$.
\end{dinhnghia}

\begin{dinhly}
  Mọi cơ sở của $V$ có cùng số phần tử. Số đó gọi là \textbf{số chiều} của $V$, ký hiệu $\dim V$.
\end{dinhly}

\subsection{Ánh xạ tuyến tính}

\begin{dinhnghia}
  Ánh xạ $T: V \to W$ gọi là \textbf{tuyến tính} nếu:
  \begin{itemize}
    \item $T(u + v) = T(u) + T(v)$,
    \item $T(\lambda u) = \lambda T(u)$.
  \end{itemize}
\end{dinhnghia}

\begin{dinhnghia}
  \textbf{Ma trận của ánh xạ tuyến tính} trong cơ sở $\mathcal{B}_V$ và $\mathcal{B}_W$: ma trận $A$ sao cho $[T(v)]_{\mathcal{B}_W} = A [v]_{\mathcal{B}_V}$.
\end{dinhnghia}

%% ================================================================
%%  CHƯƠNG 5: TRỊ RIÊNG, VECTƠ RIÊNG, DẠNG TOÀN PHƯƠNG
%% ================================================================
\section{Trị riêng và vectơ riêng}

\begin{dinhnghia}
  Số $\lambda$ gọi là \textbf{trị riêng} của ma trận $A$ nếu tồn tại vectơ $x \ne 0$ sao cho $Ax = \lambda x$. Vectơ $x$ gọi là \textbf{vectơ riêng} ứng với $\lambda$.
\end{dinhnghia}

\begin{congthuc}
  $\lambda$ là trị riêng $\Leftrightarrow$ $\det(A - \lambda I) = 0$ (phương trình đặc trưng).
\end{congthuc}

\begin{dinhnghia}
  \textbf{Dạng toàn phương} trên $\mathbb{R}^n$: $Q(x) = x^T A x = \sum_{i,j} a_{ij} x_i x_j$, với $A$ ma trận đối xứng.
\end{dinhnghia}

\begin{dinhly}[title={Dạng chính tắc}]
  Mọi dạng toàn phương đều đưa được về dạng chính tắc $Q = \lambda_1 y_1^2 + \cdots + \lambda_n y_n^2$ bằng phép đổi biến trực giao, với $\lambda_i$ là trị riêng của $A$.
\end{dinhly}

\end{document}
